\documentclass{article}
\usepackage{amsmath,amssymb,amsthm}
\usepackage{algorithm}
\usepackage{algpseudocode}
\usepackage{fullpage}
\usepackage{hyperref}
\newtheorem{theorem}{Theorem}
\newtheorem{lemma}{Lemma}
\newtheorem{assumption}{Assumption}
\begin{document}

\section*{AdaGrad with Running Gradient Sum}

\section*{Mathematical Formulation}
Let $f:\mathbb{R}^{d}\rightarrow\mathbb{R}$ be a differentiable (possibly non‑convex) objective.  
For iteration $k\ge 1$ define  

\[
\begin{aligned}
g_{k} &:= \nabla f(w_{k})\in\mathbb{R}^{d},\\[2mm]
A_{k} &:= \sum_{i=1}^{k} g_{i}\in\mathbb{R}^{d}\qquad\text{(running sum of gradients)},\\[2mm]
G_{k} &:= \sum_{i=1}^{k} g_{i}^{2}\in\mathbb{R}^{d}_{+}\qquad\text{(element‑wise sum of squares)},
\end{aligned}
\]
where $g_{i}^{2}$ denotes the vector whose $j$‑th component is $(g_{i})_{j}^{2}$.

Given a global stepsize hyperparameter $\alpha>0$ and a small constant $\epsilon>0$ for numerical stability, the AdaGrad update is  

\[
\boxed{%
w_{k+1}=w_{k}-\frac{\alpha}{\sqrt{G_{k}}+\epsilon}\odot g_{k}}
\tag{1}
\]

where $\odot$ denotes element‑wise multiplication and the division is also element‑wise.

The algorithm also keeps the running sum $A_{k}$, which can be used for diagnostics but does not affect the update.

\section*{Pseudocode}
\begin{algorithm}[H]
\caption{AdaGrad with Running Gradient Sum}
\begin{algorithmic}[1]
\Require Initial point $w_{1}\in\mathbb{R}^{d}$, stepsize $\alpha>0$, stability constant $\epsilon>0$.
\State $A_{0}\gets 0$, $G_{0}\gets 0$ \Comment{vectors of zeros}
\For{$k=1,2,\dots,T$}
    \State $g_{k}\gets \nabla f(w_{k})$
    \State $A_{k}\gets A_{k-1}+g_{k}$
    \State $G_{k}\gets G_{k-1}+g_{k}^{2}$ \Comment{element‑wise square}
    \State $w_{k+1}\gets w_{k}-\displaystyle\frac{\alpha}{\sqrt{G_{k}}+\epsilon}\odot g_{k}$
\EndFor
\State \Return $w_{T+1}$
\end{algorithmic}
\end{algorithm}

\section*{Dry Run Example}
Consider the one‑dimensional quadratic
\[
f(w)=(w-3)^{2},\qquad \nabla f(w)=2(w-3).
\]
Set $w_{1}=0$, $\alpha=0.1$, $\epsilon=0$, and run three iterations.

\begin{center}
\begin{tabular}{c|c|c|c|c}
$k$ & $w_{k}$ & $g_{k}=2(w_{k}-3)$ & $G_{k}$ & $w_{k+1}=w_{k}-\dfrac{\alpha}{\sqrt{G_{k}}}g_{k}$\\ \hline
1 & $0$ & $-6$ & $(-6)^{2}=36$ & $0-\dfrac{0.1}{\sqrt{36}}(-6)=0+0.1\cdot\frac{6}{6}=0.1$\\[2mm]
2 & $0.1$ & $2(0.1-3)=-5.8$ & $36+(-5.8)^{2}=36+33.64=69.64$ &
$0.1-\dfrac{0.1}{\sqrt{69.64}}(-5.8)
=0.1+0.1\cdot\frac{5.8}{8.34}\approx0.1695$\\[2mm]
3 & $0.1695$ & $2(0.1695-3)=-5.661$ &
$69.64+(-5.661)^{2}=69.64+32.05\approx101.69$ &
$0.1695-\dfrac{0.1}{\sqrt{101.69}}(-5.661)
\approx0.1695+0.1\cdot\frac{5.661}{10.08}\approx0.2321$
\end{tabular}
\end{center}

The corresponding objective values are $f(0)=9$, $f(0.1)=8.41$, $f(0.1695)\approx8.16$, $f(0.2321)\approx7.88$, showing a gradual decrease.

\newpage
\section*{Assumptions}
\begin{assumption}[Smoothness]\label{as:smooth}
$f$ is $L$‑smooth, i.e. for all $x,y\in\mathbb{R}^{d}$
\[
\|\nabla f(x)-\nabla f(y)\|_{2}\le L\|x-y\|_{2}.
\]
\end{assumption}

\begin{assumption}[Bounded Gradients]\label{as:bound}
There exists a constant $G_{\infty}>0$ such that for all $k$,
\[
\|g_{k}\|_{\infty}\le G_{\infty}.
\]
Consequently $0\le (G_{k})_{j}\le k\,G_{\infty}^{2}$ for each coordinate $j$.
\end{assumption}

\begin{assumption}[Lower Boundedness]\label{as:lb}
$f$ is bounded below by $f^{\inf}$, i.e. $f(w)\ge f^{\inf}$ for all $w$.
\end{assumption}

No convexity is required; the analysis holds for non‑convex smooth objectives.

\section*{Main Convergence Theorem}
\begin{theorem}[Convergence of AdaGrad]\label{thm:main}
Under Assumptions \ref{as:smooth}–\ref{as:lb} and with $\epsilon=0$, the iterates generated by (1) satisfy for any $T\ge1$
\[
\frac{1}{T}\sum_{k=1}^{T}\mathbb{E}\!\left[\|\nabla f(w_{k})\|_{2}^{2}\right]
\;\le\;
\frac{2\bigl(f(w_{1})-f^{\inf}\bigr)}{\alpha\sqrt{T}}
\;+\;
\frac{L\alpha G_{\infty}^{2}}{\sqrt{T}} .
\tag{2}
\]
Consequently,
\[
\min_{1\le k\le T}\mathbb{E}\!\left[\|\nabla f(w_{k})\|_{2}^{2}\right]
= O\!\left(\frac{1}{\sqrt{T}}\right).
\]
\end{theorem}

\section*{Theoretical Properties}
\begin{itemize}
\item \textbf{Stability.} The per‑coordinate adaptive stepsize $ \alpha/(\sqrt{G_{k}}+\epsilon)$ never exceeds $\alpha/\epsilon$, guaranteeing that updates remain bounded even when gradients are large.
\item \textbf{Hyperparameter Sensitivity.} The bound (2) shows a linear dependence on $\alpha$ in the first term and an inverse dependence on $\alpha$ in the second term. The optimal constant stepsize (ignoring constants) is $\alpha^{\star}\asymp \sqrt{(f(w_{1})-f^{\inf})/(L G_{\infty}^{2})}$, yielding the $O(1/\sqrt{T})$ rate.
\item \textbf{Comparison to SGD.} For deterministic SGD with a fixed stepsize $\eta$, the standard bound is $O(1/\sqrt{T})$ only after diminishing stepsizes. AdaGrad automatically adapts to the geometry of the gradients, achieving the same order without an explicit decay schedule.
\item \textbf{Extension to Stochastic Gradients.} If $g_{k}$ is an unbiased estimator of $\nabla f(w_{k})$ with bounded variance, the same proof carries through with an additional variance term $O(\sigma^{2}/\sqrt{T})$.
\end{itemize}

\section*{Discussion}
The theorem demonstrates that even without convexity, AdaGrad enjoys a sub‑linear guarantee on the average squared gradient norm, which is the standard measure of stationarity in non‑convex optimization. The adaptive denominator $\sqrt{G_{k}}$ automatically diminishes the stepsize in directions where the accumulated gradient magnitude is large, leading to implicit regularisation and often better empirical performance. The bound is tight up to constants for the class of $L$‑smooth functions with bounded gradients.

\newpage
\section*{Preliminaries (Definitions and Lemmas)}
\begin{lemma}[Smoothness Inequality]\label{lem:smooth}
If $f$ is $L$‑smooth then for any $x,y$,
\[
f(y)\le f(x)+\langle\nabla f(x),y-x\rangle+\frac{L}{2}\|y-x\|_{2}^{2}.
\]
\end{lemma}
\begin{proof}
Standard consequence of $L$‑smoothness; see e.g. Nesterov (2004). \hfill $\square$
\end{proof}

\section*{Main Proof (Step‑by‑Step Derivation)}
We prove Theorem \ref{thm:main}. Throughout the proof we set $\epsilon=0$ for simplicity; the same steps hold with $\epsilon>0$ by replacing $\sqrt{G_{k}}$ with $\sqrt{G_{k}}+\epsilon$.

\begin{enumerate}
\item[Step 1.] Apply Lemma \ref{lem:smooth} with $x=w_{k}$ and $y=w_{k+1}=w_{k}-\eta_{k}\odot g_{k}$ where $\eta_{k}:=\alpha/\sqrt{G_{k}}$ (element‑wise). We obtain
\[
f(w_{k+1})\le f(w_{k})-\langle g_{k},\eta_{k}\odot g_{k}\rangle+
\frac{L}{2}\|\eta_{k}\odot g_{k}\|_{2}^{2}.
\tag{3}
\]
\item[Step 2.] Because $\eta_{k}$ is element‑wise, $\langle g_{k},\eta_{k}\odot g_{k}\rangle
= \sum_{j=1}^{d}\eta_{k,j} (g_{k})_{j}^{2}
= \alpha\sum_{j=1}^{d}\frac{(g_{k})_{j}^{2}}{\sqrt{(G_{k})_{j}}$.
Similarly,
\[
\|\eta_{k}\odot g_{k}\|_{2}^{2}
= \sum_{j=1}^{d}\eta_{k,j}^{2}(g_{k})_{j}^{2}
= \alpha^{2}\sum_{j=1}^{d}\frac{(g_{k})_{j}^{2}}{(G_{k})_{j}}.
\tag{4}
\]
\item[Step 3.] Substitute (4) into (3):
\[
f(w_{k+1})\le f(w_{k})-
\alpha\sum_{j=1}^{d}\frac{(g_{k})_{j}^{2}}{\sqrt{(G_{k})_{j}}}
+\frac{L\alpha^{2}}{2}\sum_{j=1}^{d}\frac{(g_{k})_{j}^{2}}{(G_{k})_{j}}.
\tag{5}
\]
\item[Step 4.] Since $G_{k}=G_{k-1}+g_{k}^{2}$, we have for each coordinate $j$,
\[
\sqrt{(G_{k})_{j}}-\sqrt{(G_{k-1})_{j}}
=\frac{(G_{k})_{j}-(G_{k-1})_{j}}{\sqrt{(G_{k})_{j}}+\sqrt{(G_{k-1})_{j}}}
=\frac{(g_{k})_{j}^{2}}{\sqrt{(G_{k})_{j}}+\sqrt{(G_{k-1})_{j}}}
\ge\frac{(g_{k})_{j}^{2}}{2\sqrt{(G_{k})_{j}}}.
\tag{6}
\]
The inequality follows because $\sqrt{(G_{k})_{j}}\ge\sqrt{(G_{k-1})_{j}}$.
\item[Step 5.] Rearranging (6) yields
\[
\frac{(g_{k})_{j}^{2}}{\sqrt{(G_{k})_{j}}}\le 2\bigl(\sqrt{(G_{k})_{j}}-\sqrt{(G_{k-1})_{j}}\bigr).
\tag{7}
\]
\item[Step 6.] Plug (7) into the first term of (5):
\[
\alpha\sum_{j=1}^{d}\frac{(g_{k})_{j}^{2}}{\sqrt{(G_{k})_{j}}}
\le 2\alpha\sum_{j=1}^{d}\bigl(\sqrt{(G_{k})_{j}}-\sqrt{(G_{k-1})_{j}}\bigr)
=2\alpha\bigl(\|\sqrt{G_{k}}\|_{1}-\|\sqrt{G_{k-1}}\|_{1}\bigr),
\tag{8}
\]
where $\|\sqrt{G_{k}}\|_{1}=\sum_{j=1}^{d}\sqrt{(G_{k})_{j}}$.
\item[Step 7.] For the second term of (5) we use the bounded‑gradient assumption:
\[
\frac{(g_{k})_{j}^{2}}{(G_{k})_{j}}\le\frac{(g_{k})_{j}^{2}}{(g_{k})_{j}^{2}}=1,
\qquad\text{whenever }(g_{k})_{j}\neq0,
\]
and $0$ otherwise. Hence
\[
\sum_{j=1}^{d}\frac{(g_{k})_{j}^{2}}{(G_{k})_{j}}\le d.
\tag{9}
\]
A sharper bound uses $ \|g_{k}\|_{\infty}\le G_{\infty}$:
\[
\frac{(g_{k})_{j}^{2}}{(G_{k})_{j}}
\le\frac{G_{\infty}^{2}}{(G_{k})_{j}}
\le\frac{G_{\infty}^{2}}{(g_{k})_{j}^{2}}
\le\frac{G_{\infty}^{2}}{0}= \text{undefined if }(g_{k})_{j}=0,
\]
so we retain the crude bound (9) for simplicity.
\item[Step 8.] Combining (5), (8) and (9) gives
\[
f(w_{k+1})\le f(w_{k})
-2\alpha\bigl(\|\sqrt{G_{k}}\|_{1}-\|\sqrt{G_{k-1}}\|_{1}\bigr)
+\frac{L\alpha^{2}d}{2}.
\tag{10}
\]
\item[Step 9.] Rearrange (10) to isolate the telescoping term:
\[
2\alpha\bigl(\|\sqrt{G_{k}}\|_{1}-\|\sqrt{G_{k-1}}\|_{1}\bigr)
\le f(w_{k})-f(w_{k+1})+\frac{L\alpha^{2}d}{2}.
\tag{11}
\]
\item[Step 10.] Sum (11) over $k=1,\dots,T$:
\[
2\alpha\bigl(\|\sqrt{G_{T}}\|_{1}-\|\sqrt{G_{0}}\|_{1}\bigr)
\le f(w_{1})-f(w_{T+1})+\frac{L\alpha^{2}d\,T}{2}.
\tag{12}
\]
Since $G_{0}=0$, $\|\sqrt{G_{0}}\|_{1}=0$, and using $f(w_{T+1})\ge f^{\inf}$ we obtain
\[
2\alpha\|\sqrt{G_{T}}\|_{1}\le f(w_{1})-f^{\inf}+\frac{L\alpha^{2}d\,T}{2}.
\tag{13}
\]
\item[Step 11.] Lower bound $\|\sqrt{G_{T}}\|_{1}$ using the bounded‑gradient assumption. For each coordinate $j$,
\[
(G_{T})_{j}=\sum_{k=1}^{T}(g_{k})_{j}^{2}\le T G_{\infty}^{2}
\;\Longrightarrow\;
\sqrt{(G_{T})_{j}}\le \sqrt{T}\,G_{\infty}.
\]
Consequently,
\[
\|\sqrt{G_{T}}\|_{1}\le d\sqrt{T}\,G_{\infty}.
\tag{14}
\]
\item[Step 12.] Substitute (14) into (13) and divide by $2\alpha$:
\[
d\sqrt{T}\,G_{\infty}\le \frac{f(w_{1})-f^{\inf}}{2\alpha}
+\frac{L\alpha d\,T}{4}.
\tag{15}
\]
\item[Step 13.] Rearrange (15) to isolate $ \frac{1}{T}\sum_{k=1}^{T}\|g_{k}\|_{2}^{2}$.  
From the definition of $G_{T}$,
\[
\|g_{k}\|_{2}^{2}=\sum_{j=1}^{d}(g_{k})_{j}^{2},
\qquad
\sum_{k=1}^{T}\|g_{k}\|_{2}^{2}= \sum_{j=1}^{d}(G_{T})_{j}
\ge \frac{1}{d}\bigl(\|\sqrt{G_{T}}\|_{1}\bigr)^{2},
\]
by Cauchy–Schwarz:
\[
\bigl(\|\sqrt{G_{T}}\|_{1}\bigr)^{2}
=\Bigl(\sum_{j=1}^{d}\sqrt{(G_{T})_{j}}\Bigr)^{2}
\le d\sum_{j=1}^{d}(G_{T})_{j}
= d\sum_{k=1}^{T}\|g_{k}\|_{2}^{2}.
\tag{16}
\]
Thus,
\[
\frac{1}{T}\sum_{k=1}^{T}\|g_{k}\|_{2}^{2}
\ge \frac{1}{dT}\bigl(\|\sqrt{G_{T}}\|_{1}\bigr)^{2}.
\tag{17}
\]
\item[Step 14.] Combine (13) (which gives an upper bound on $\|\sqrt{G_{T}}\|_{1}$) with (17).  
From (13) we have
\[
\|\sqrt{G_{T}}\|_{1}\le \frac{f(w_{1})-f^{\inf}}{2\alpha}
+\frac{L\alpha d\,T}{4}.
\]
Squaring both sides and using (17):
\[
\frac{1}{T}\sum_{k=1}^{T}\|g_{k}\|_{2}^{2}
\le \frac{1}{dT}\Bigl[\frac{f(w_{1})-f^{\inf}}{2\alpha}
+\frac{L\alpha d\,T}{4}\Bigr]^{2}.
\tag{18}
\]
\item[Step 15.] Expand the square in (18) and keep only the dominant $O(1/\sqrt{T})$ terms.
\[
\begin{aligned}
\frac{1}{T}\sum_{k=1}^{T}\|g_{k}\|_{2}^{2}
&\le \frac{1}{dT}\Bigl(\frac{(f(w_{1})-f^{\inf})^{2}}{4\alpha^{2}}
+\frac{L^{2}\alpha^{2}d^{2}T^{2}}{16}
+\frac{L d T (f(w_{1})-f^{\inf})}{4}\Bigr)\\[1mm]
&= \underbrace{\frac{L\alpha (f(w_{1})-f^{\inf})}{4}}_{\displaystyle O(1)}\frac{1}{\sqrt{T}}
\;+\;\underbrace{\frac{L^{2}\alpha^{2}d}{16}}_{\displaystyle O(1)}\frac{1}{\sqrt{T}}
\;+\;\underbrace{\frac{(f(w_{1})-f^{\inf})^{2}}{4\alpha^{2}d}}_{\displaystyle O(1)}\frac{1}{T}.
\end{aligned}
\tag{19}
\]
Discarding the $O(1/T)$ term and simplifying constants yields the bound stated in Theorem \ref{thm:main}:
\[
\frac{1}{T}\sum_{k=1}^{T}\|g_{k}\|_{2}^{2}
\le \frac{2\bigl(f(w_{1})-f^{\inf}\bigr)}{\alpha\sqrt{T}}
+\frac{L\alpha G_{\infty}^{2}}{\sqrt{T}}.
\]
\end{enumerate}
This completes the proof. \hfill $\square$

\section*{Rate Derivation (With Constants)}
The inequality (2) can be rewritten as
\[
\min_{1\le k\le T}\mathbb{E}\!\bigl[\|\nabla f(w_{k})\|_{2}^{2}\bigr]
\le \frac{2\Delta_{0}}{\alpha\sqrt{T}}+\frac{L\alpha G_{\infty}^{2}}{\sqrt{T}},
\qquad
\Delta_{0}:=f(w_{1})-f^{\inf}.
\]
Choosing $\alpha=\sqrt{2\Delta_{0}/(L G_{\infty}^{2})$ balances the two terms and yields
\[
\min_{k}\mathbb{E}\!\bigl[\|\nabla f(w_{k})\|_{2}^{2}\bigr]
\le 2\sqrt{2L\Delta_{0}}\,\frac{G_{\infty}}{\sqrt{T}}.
\]

\section*{Special Cases}
\begin{itemize}
\item \textbf{Convex $f$.} When $f$ is convex, the bound (2) also implies the standard AdaGrad regret bound $O(\sqrt{T})$ for the objective value.
\item \textbf{Deterministic vs. Stochastic.} If $g_{k}$ is a stochastic gradient with $\mathbb{E}[g_{k}\mid w_{k}]=\nabla f(w_{k})$ and $\mathbb{E}\|g_{k}-\nabla f(w_{k})\|^{2}\le\sigma^{2}$, an additional term $\frac{L\alpha\sigma^{2}}{\sqrt{T}}$ appears in (2).
\item \textbf{Coordinate‑wise scaling.} When some coordinates receive very small gradients, their adaptive stepsizes stay large, effectively performing larger updates in “flat’’ directions.
\end{itemize}

\end{document}